% !TeX root = ../../Book1.tex
% 纲要标题：### 4.1 作用、轨道与稳定子群
% 纲要位置：计划纲要.md，第 217 行。

\section{作用、轨道与稳定子群}
\label{b1:sec:0401}

正方形的对称可以置换顶点，也可以置换边或对角线。同一个群因所选集合不同而给出不同的信息，群作用正是描述这种联系的语言。考察一点能被移到哪些位置，以及哪些变换使它不动，就得到轨道和稳定子群；两者将把置换的具体行为与子群的指数联系起来。以下一律采用左作用，置换的复合从右向左进行。

% 本节正文按下列原纲要要点撰写。

% 纲要内容（仅作写作计划，尚非教材正文）：
%
% * 群作用与到置换群的同态
% * 轨道、稳定子群与作用核；忠实性与自由性
% * 传递作用与陪集空间
%

\subsection{群作用与到置换群的同态}
\label{b1:subsec:0401:01}

给每个群元素 \(g\) 指定集合 \(X\) 的一个置换，还需要这些置换彼此相容：单位元素应当对应恒等置换，群中的乘积应当对应置换的复合。群作用把这两个要求写成每一点都满足的等式。

\begin{definition}\label{b1:def:group-action}
设 \(G\) 为群，\(X\) 为集合，并给定映射
\[
G\times X\longrightarrow X,\qquad (g,x)\longmapsto gx,
\]
如果对任意 \(g,h\in G\) 和 \(x\in X\)，都有 \(1_Gx=x\) 及 \((gh)x=g(hx)\)，则将这个映射称为 \(G\) 在 \(X\) 上的\term[qun-zuoyong]{群作用}[group action]，并将带有这个指定作用的集合 \(X\) 称为 \(G\)-集合。
\end{definition}

记 \(\lambda_g(x)=gx\)。由作用公理，\(\lambda_{g^{-1}}\) 是 \(\lambda_g\) 的逆映射，故每个 \(\lambda_g\) 都是置换。左作用中的 \((gh)x=g(hx)\) 表示先作用 \(h\)，再作用 \(g\)，因而
\[
\lambda_{gh}=\lambda_g\circ\lambda_h.
\]
因此作用等价于同态 \(G\rightarrow S_X\)。反过来，给定这样的同态并令 \(gx=\lambda_g(x)\)，同态保持单位和乘法就给出两条作用公理。这一对应也解释了为什么作用中不必另行假设每个 \(g\) 都给出双射。不过，这个同态不一定单射：不同的群元素可能诱导相同的置换。诱导恒等置换的元素构成这个同态的核，稍后称为作用的核；若 \(\lambda_g=\lambda_h\)，则落入核的是 \(h^{-1}g\)，不要求 \(g,h\) 本身都在核中。

\ProseParagraphBreak
例如，\(S_n\) 在 \(\{1,\ldots,n\}\) 上按定义作用。若 \(X\) 是这个集合的所有二元子集，则 \(\sigma\{i,j\}=\bigl\{\sigma(i),\sigma(j)\bigr\}\) 又定义一个作用。后一作用的点是无序对，不能把它与有序对作用混同：交换 \(i,j\) 的置换固定 \(\{i,j\}\)，却不固定 \((i,j)\)。

\subsection{轨道、稳定子群与作用核}
\label{b1:subsec:0401:02}

\begin{definition}\label{b1:def:orbit-stabilizer}
设群 \(G\) 作用在集合 \(X\) 上，并取 \(x\in X\)。令
\[
Gx\coloneqq\{\,gx\mid g\in G\,\},\qquad
G_x\coloneqq\{\,g\in G\mid gx=x\,\}
\]
分别称 \(Gx\) 与 \(G_x\) 为 \(x\) 的\term[guidao]{轨道}[orbit]与\term[wendingziqun]{稳定子群}[stabilizer]。
\symidx{\(Gx\)：群作用的轨道}
\symidx{\(G_x\)：点的稳定子群}
\end{definition}

这两个记号表示不同集合中的对象：\(Gx\subseteq X\) 是由点组成的轨道，回答“\(x\) 能被移到哪里”；\(G_x\subseteq G\) 是由群元素组成的稳定子群，回答“哪些元素使 \(x\) 留在原处”。

\ProseParagraphBreak
稳定子群也称稳定化子。\footnote{“稳定子群”见欧阳毅等《代数学 II：近世代数》\cite[第 2.2.1 节]{OuyangYeChen2017AlgebraII}；“稳定化子”见石生明《近世代数初步》\cite[第一章第 6 节]{Shi2007IntroModernAlgebra}。}\termidx[wendinghuazi]{稳定化子}它由所有固定 \(x\) 的群元素组成；讨论单个群元素时，直接说“固定 \(x\) 的元素 \(g\)”。

\ProseParagraphBreak
关系“\(y=gx\) 对某个 \(g\) 成立”是等价关系：单位、逆元、乘法分别给出自反性、对称性和传递性。因此各轨道两两不交并覆盖 \(X\)。稳定子群是子群，因为 \(1\in G_x\)，而 \(g,h\in G_x\) 时 \(gh^{-1}x=x\)。

\ProseParagraphBreak
要把轨道中的位置与群中的元素联系起来，需要知道哪些元素把 \(x\) 移到同一点。对 \(g,h\in G\)，有
\[
gx=hx
\iff h^{-1}g\in G_x
\iff gG_x=hG_x.
\]
也就是说，\(g\) 与 \(h\) 给出同一个位置，恰好意味着 \(g\) 可由 \(h\) 右乘一个固定 \(x\) 的元素得到。因此，把 \(x\) 移到某个指定位置的所有群元素组成一个左陪集；不同位置对应不同陪集。

\begin{theorem}[轨道—稳定子群]\label{b1:thm:orbit-stabilizer}
设群 \(G\) 作用在集合 \(X\) 上，并取 \(x\in X\)。这里 \(G/G_x\) 只表示左陪集集合，不要求 \(G_x\) 正规。映射
\[
G/G_x\longrightarrow Gx,\qquad gG_x\longmapsto gx
\]
是双射。因此有基数等式
\[
|Gx|=[G:G_x].
\]
若 \(G\) 有限，还可写成有限整数的乘积公式
\[
|G|=|Gx|\cdot|G_x|.
\]
\end{theorem}
\begin{proof}
若 \(gG_x=hG_x\)，则 \(h^{-1}g\in G_x\)，故 \(gx=hx\)，映射良定义。反之，\(gx=hx\) 推出 \(h^{-1}gx=x\)，所以 \(gG_x=hG_x\)，映射单射。轨道的定义给出满射，基数等式随即成立。

有限情形中，\cref{b1:thm:lagrange}用于子群 \(G_x\leq G\)，给出 \(|G|=[G:G_x]|G_x|\)；再用基数等式代入，就得到乘积公式。
\end{proof}

这个结果称为\term[guidao-wendingziqun-dingli]{轨道—稳定子群定理}[orbit--stabilizer theorem]。它把“能把 \(x\) 移到多少位置”与“有多少变换保持 \(x\) 不动”联系起来。在 \(S_4\) 对二元子集的作用中，\(\{1,2\}\) 的稳定子群可以独立交换 \(1,2\) 和 \(3,4\)，恰有 \(4\) 个元素；轨道大小为 \(24/4=6\)，与六个二元子集相符。

\ProseParagraphBreak
若把同一轨道中的基点从 \(x\) 换成 \(y=gx\)，则 \(G_y=gG_xg^{-1}\)：元素 \(h\) 固定 \(gx\) 当且仅当 \(g^{-1}hg\) 固定 \(x\)。因此，换基点时，描述该点的稳定子群相应地通过共轭改变；同一轨道内的稳定子群通常不相等，但总互相共轭，特别是具有相同阶。

\ProseParagraphBreak
稳定子群考察哪些群元素固定一个指定点。若要判断一个作用能否区分全部群元素，则要同时考察所有点：\(g,h\in G\) 给出相同置换，当且仅当 \(h^{-1}g\) 固定每一点。

\begin{definition}\label{b1:def:faithful-action}
设群 \(G\) 作用在集合 \(X\) 上，对应同态为 \(\rho:G\rightarrow S_X\)。使 \(X\) 中每一点都不动的群元素组成作用的核，记为
\[
K\coloneqq\ker\rho
=\{\,g\in G\mid gx=x\;\text{对所有}\;x\in X\,\}.
\]
如果只有单位元素固定全部的点，即 \(K=\{1_G\}\)，则将此作用称为\term[zhongshi-zuoyong]{忠实作用}[faithful action]。如果对任意 \(x\in X\) 和 \(g\in G\)，等式 \(gx=x\) 都蕴含 \(g=1_G\)，则将此作用称为\term[ziyou-zuoyong]{自由作用}[free action]。
\end{definition}

按定义，固定指定点 \(x\) 的元素组成 \(G_x\)，同时固定每一点的元素则组成这些稳定子群的交，因此
\[
\ker\rho=\bigcap_{x\in X}G_x.
\]
作用的核是同态的核，因而是 \(G\) 的正规子群。由\cref{b1:prop:kernel-fibers}，忠实等价于 \(\rho\) 单射；自由则等价于每个稳定子群都平凡。\footnote{当 \(X=\varnothing\) 时，空族子群之交约定为整个群 \(G\)。空集只有一个置换，也只有一个 \(G\)-作用，其核为 \(G\)。自由性的全称条件此时自动满足，但只有平凡群的作用忠实。}

\ProseParagraphBreak
若 \(X\ne\varnothing\)，自由作用一定忠实：固定全部点的元素也固定任取的一点，所以只能是单位。反过来不成立，\(S_3\) 在三个点上的自然作用忠实，但每个换位仍固定一个点。忠实性要求一个元素只要固定所有点就必须为单位；自由性要求一个元素只要固定任一点就必须为单位。

\subsection{传递作用与陪集空间}
\label{b1:subsec:0401:03}

\begin{definition}\label{b1:def:transitive-action}
设群 \(G\) 作用在非空集合 \(X\) 上。如果对任意 \(x,y\in X\)，都存在 \(g\in G\) 使 \(gx=y\)，则称此作用为\term[chuandi-zuoyong]{传递作用}[transitive action]。

对任意两个 \(G\)-集合 \(X,Y\)，如果映射 \(f:X\rightarrow Y\) 对所有 \(g\in G\)、\(x\in X\) 都满足 \(f(gx)=gf(x)\)，则将 \(f\) 称为\term[dengbian-yingshe]{等变映射}[equivariant map]。如果等变映射还是双射，则将它称为 \(G\)-集合的\term[dengbian-tonggou]{等变同构}[equivariant isomorphism]。
\end{definition}

传递性等价于 \(X\) 只有一个轨道。对非空的 \(G\)-集合 \(X\)，这三种性质可比较如下。

\begin{center}
\begin{tabularx}{\linewidth}{@{}lXX@{}}
\toprule
性质 & 条件及量词 & 轨道或稳定子群的描述 \\
\midrule
忠实 & 若 \(g\in G\) 对所有 \(x\in X\) 都满足 \(gx=x\)，则 \(g=1_G\) & \(\displaystyle\bigcap_{x\in X}G_x=\{1_G\}\) \\
自由 & 对每个 \(x\in X\)，若 \(g\in G\) 满足 \(gx=x\)，则 \(g=1_G\) & 对每个 \(x\in X\)，都有 \(G_x=\{1_G\}\) \\
传递 & 对任意 \(x,y\in X\)，存在 \(g\in G\) 使 \(gx=y\) & 对每个 \(x\in X\)，都有 \(Gx=X\)，即只有一个轨道 \\
\bottomrule
\end{tabularx}
\end{center}

自由要求每个稳定子群都平凡，忠实只要求它们的交平凡，所以非空情形下自由蕴含忠实。传递性则考察轨道数，与一个元素能固定多少点是不同的条件。

\ProseParagraphBreak
等变性说明先作用再映射与先映射再作用得到相同结果；定义中的两个作用可以不同，但使用的是同一个群 \(G\)。例如，让 \(S_3\) 作用在集合
\[
X=\{\,(i,j)\mid 1\le i,j\le 3,\ i\ne j\,\},\qquad
Y=\{\,\{i,j\}\mid 1\le i<j\le 3\,\}
\]
上，分别按 \(\sigma(i,j)=(\sigma(i),\sigma(j))\) 和 \(\sigma\{i,j\}=\{\sigma(i),\sigma(j)\}\) 作用。忘掉次序的映射
\[
f:X\longrightarrow Y,\qquad (i,j)\longmapsto\{i,j\}
\]
满足
\[
f\bigl(\sigma(i,j)\bigr)
=\{\sigma(i),\sigma(j)\}
=\sigma f(i,j),
\]
因而等变。它把 \((1,2)\) 与 \((2,1)\) 映到同一点，所以等变映射不必单射；等变同构还要求映射为双射。

\ProseParagraphBreak
等变双射 \(F:X\rightarrow Y\) 的逆映射也等变。事实上，若 \(y=F(x)\)，则
\[
F^{-1}(gy)=F^{-1}\bigl(gF(x)\bigr)
=F^{-1}\bigl(F(gx)\bigr)
=gx=gF^{-1}(y).
\]
这里比较的是集合及其指定作用；这些集合本身不必是群，等变同构也不是群同构的另一名称。

\ProseParagraphBreak
任取子群 \(H\leq G\)，公式 \(a(gH)=(ag)H\) 给出 \(G\) 在左陪集集合 \(G/H\) 上的传递作用。若 \(gH=g'H\)，则 \(agH=ag'H\)，所以公式与陪集代表元的选择无关。点 \(H\) 的稳定子群恰为 \(H\)。

\ProseParagraphBreak
重建一个传递作用时，可以先选定基点 \(x\)，再记录子群 \(G_x\)。轨道—稳定子群定理给出的双射说明，轨道中的点可由 \(G_x\) 的左陪集描述；还需核对这个双射是否保留作用。换基点会使稳定子群共轭，因此忘掉基点后，相应的数据应当是子群的共轭类。下面的命题把这两点合在一起。

\begin{proposition}\label{b1:prop:transitive-coset}
设群 \(G\) 传递地作用在非空集合 \(X\) 上。对任意 \(x\in X\)，\(G\)-集合 \(X\) 都等变同构于左陪集集合 \(G/G_x\)。此外，对任意子群 \(H,K\leq G\)，两个左陪集作用 \(G/H\)、\(G/K\) 等变同构，当且仅当 \(H,K\) 在 \(G\) 中共轭。
\end{proposition}
\begin{proof}
由\cref{b1:thm:orbit-stabilizer}，映射 \(f:gG_x\mapsto gx\) 是从 \(G/G_x\) 到轨道 \(Gx\) 的双射；传递性使 \(Gx=X\)。又有 \(f(agG_x)=agx=af(gG_x)\)，所以此双射等变。

若 \(F:G/H\rightarrow G/K\) 是等变双射，写 \(F(H)=aK\)。对任意 \(u\in G\)，等变性给出 \(F(uH)=uF(H)\)，再由 \(F\) 单射，得到 \(uH=H\) 当且仅当 \(uaK=aK\)。两边分别等价于 \(u\in H\) 和 \(u\in aKa^{-1}\)，因此 \(H=aKa^{-1}\)。

反过来，若 \(H=aKa^{-1}\)，则 \(gH\mapsto gaK\) 良定义：把 \(g\) 改为 \(gh\)、\(h\in H\) 时，由 \(a^{-1}ha\in K\) 得到 \(haK=aK\)。同理由 \(K=a^{-1}Ha\) 可知 \(gK\mapsto ga^{-1}H\) 良定义；两者互为逆映射，且都与左乘作用相容，所以给出等变同构。
\end{proof}

一般作用按轨道分解为这样的传递部分。这里不应把“传递”理解为“每个非单位元素都移动每一点”：例如 \(S_3\) 的自然作用传递，但换位固定第三个点。
